\section{Uniqueness of the continuum law}\label{sec:uniqueness}

The preceding sections identify every bounded-domain subsequential
limit with Equation~\eqref{mod:SDE}. We now prove that the continuum
prescription determines one law. The argument uses weak uniqueness
for a measurable drift functional; it requires no Lipschitz estimate
for the dependence of the boundary value problem on a slit.

\subsection{The coefficient as a path functional}

\begin{lemma}\label{uni:measurable}
There is a fixed nonanticipating progressively measurable extension
of $C_t$ to continuous driving paths, agreeing with
Equation~\eqref{mod:C} whenever its integral is finite. The same
extension can be used for every candidate law in
Definition~\ref{def:law}.
\end{lemma}

\begin{proof}
For a continuous driver and a fixed time, the Loewner ordinary
differential equation determines the maps on all unswallowed
compact sets. The inverse maps depend continuously on the driver
on such protected sets. This also gives the marked kernel topology
of $O_t$: in the centered chart the two marks are always
$0,\infty$.

Proposition~\ref{mod:bvp-unique} supplies a unique $h_t$ and
kernel-continuous interior values and gradients. To spell out the
measurability consequence, exhaust $\Hh$ by deterministic compact
sets $K_j$ avoiding $0$ and the real line. The integrals
\[
 C_t^{(j)}=\frac4\pi\int_{K_j}P(v)m|\psi_t'(v)|
                  |\grad(h_t\circ\psi_t)(v)|\,\dd A(v)
\]
are Borel functions of the stopped driver. One can first use compact
sets on which the maps and their derivatives satisfy fixed
protection bounds and then exhaust these bounds. The continuity
assertion in Proposition~\ref{mod:bvp-unique} applies because all
physical active domains lie in the fixed bounded domain $\Omega$.
Taking an increasing exhaustion $K_j$ makes
$C_t=\lim_j C_t^{(j)}$ an extended nonnegative Borel function.

All these constructions at time $t$ use only the path up to $t$.
More explicitly they are Borel functions of $(t,x^t)$, where
$x^t(s)=x(s\wedge t)$ and the stopped-path map is continuous on
each finite time interval. They are therefore progressively
measurable. Set the extension to zero on the measurable set where
the extended integral is infinite or the prescription is undefined.
The admissibility conditions retain the requirement that the actual
integral be finite almost everywhere, so this convention adds no
solutions by substituting a value on a positive-measure set of
undefined coefficients.
\end{proof}

\subsection{A finite-energy uniqueness lemma}

The following formulation isolates a small point about energy
stopping. Energy on an arbitrary reference Brownian path is allowed
to become infinite immediately to the right of a time. Its first
level need not then be attained continuously. The proof uses
attainment under the candidate solution to fix the stopping rule
under the comparison Wiener law.

\begin{lemma}\label{uni:weak}
Let $\kappa>0$, and let $b(t,x)$ be a fixed progressively measurable
nonanticipating real functional on continuous paths starting at zero.
There is at most one law of a continuous weak solution of
\begin{equation}\label{uni:generic-SDE}
 \dd X_t=\sqrt\kappa\,\dd B_t+b(t,X)\dd t,
 \qquad X_0=0,
\end{equation}
among solutions satisfying
\begin{equation}\label{uni:generic-energy}
 \int_0^T b(t,X)^2\dd t<\infty\quad\hbox{almost surely}
\end{equation}
for every finite $T$.
\end{lemma}

\begin{proof}
Fix a finite horizon $T$. On canonical path space define extended
energy and its first level by
\[
 A_t(x)=\int_0^t b(s,x)^2\dd s\in[0,\infty],\qquad
 \sigma_n(x)=\inf\{t\ge0:A_t(x)\ge n\}\wedge T.
\]
Use the usual right-continuous augmentation for stopping times.
Even for extended energy, $\sigma_n$ is a stopping time: the
strict sublevel events $\{\sigma_n<t\}$ can be expressed using
rational times less than $t$, and right continuity gives the weak
sublevel events.

Let $X$ be any candidate solution and put $\tau_n=\sigma_n(X)$.
By Equation~\eqref{uni:generic-energy}, $A_t(X)$ is finite and
continuous on $[0,T]$, and $A_{\tau_n}(X)\le n$. The stochastic
exponential
\begin{equation}\label{uni:tilt}
 Z_n=\exp\left\{-\frac1{\sqrt\kappa}
       \int_0^{\tau_n} b(s,X)\dd B_s
       -\frac1{2\kappa}A_{\tau_n}(X)\right\}
\end{equation}
is a true martingale density by the bounded-energy Novikov
criterion; for this criterion and the following Girsanov change of
measure, see \cite[Section~5.3, Theorems~5.3.1--5.3.2 and (5.10)]{Lawler}.
It is strictly positive. Under the equivalent measure
$Q_n$ with density $Z_n$, the process
\[
 \widetilde B_t=B_t+\frac1{\sqrt\kappa}
                  \int_0^{t\wedge\tau_n}b(s,X)\dd s,
 \qquad 0\le t\le T,
\]
is Brownian motion. Set $Y_t=\sqrt\kappa\widetilde B_t$.
It has the full Wiener law $\mathsf W_\kappa$ on $[0,T]$ and
agrees with $X$ through $\tau_n$.

Nonanticipation implies $A_t(Y)=A_t(X)$ for $t\le\tau_n$.
If $\tau_n<T$, that energy is strictly below $n$ earlier and
equals $n$ at $\tau_n$. Consequently
\begin{equation}\label{uni:same-stop}
 \sigma_n(Y)=\tau_n.
\end{equation}
If $\tau_n=T$, the paths agree on the whole horizon and the same
identity holds. Possible nonintegrability of $b(s,Y)^2$ immediately
after $\tau_n$ cannot change this first level: it has already been
attained. This observation also proves, under the fixed full Wiener
law, that
$A_{\sigma_n}(Y)\le n$ almost surely. A hypothetical Wiener path
whose energy jumps from below $n$ to infinity before this first
attained prefix has Wiener measure zero whenever a finite-energy
solution exists.

The inverse of Equation~\eqref{uni:tilt}, written under $Q_n$, is
the canonical Wiener functional
\begin{equation}\label{uni:inverse}
 R_n(y)=\exp\left\{\frac1\kappa
       \int_0^{\sigma_n(y)} b(s,y)\dd y_s
       -\frac1{2\kappa}A_{\sigma_n(y)}(y)\right\}.
\end{equation}
The stopped energy bound just established defines this stochastic
integral in $L^2(\mathsf W_\kappa)$, and hence as a measurable
functional of the canonical path, up to a single Wiener-null set.
One may use predictable simple-process approximations to choose
that version once for all candidate solutions. In particular, both
the stopped path and its inverse density are fixed by the same
Wiener law. For every bounded Borel functional $F$,
\begin{equation}\label{uni:stopped-law}
 \E F(X^{\tau_n})
 =\int F(y^{\sigma_n(y)})R_n(y)\,
                         \mathsf W_\kappa(\dd y).
\end{equation}
The right side does not depend on the candidate solution. Thus all
candidate stopped laws coincide.

Under each candidate, $A_T(X)<\infty$ almost surely; hence
$\tau_n=T$ for all sufficiently large $n$ on each sample path.
Bounded convergence in Equation~\eqref{uni:stopped-law} proves
equality of the full laws on $[0,T]$. No assertion that
$\sigma_n\to T$ under unweighted Wiener measure is needed. Since
$T$ was arbitrary, the laws on the full continuous path space
coincide.
\end{proof}

\subsection{Existence, curves and bounded approximations}

\begin{proposition}\label{uni:bounded}
Definition~\ref{def:law} determines exactly one law. Every
admissible approximation whose physical domains have a common
diameter bound converges to this law in the disc curve topology.
\end{proposition}

\begin{proof}
Consider first any such approximating sequence. The marked
basepoint normalization and common diameter bound place all its
domains in a fixed physical ball. Proposition~\ref{geo:compactness}
gives tightness jointly for curves and drivers; every subsequential
limit has the full continuous Loewner trace and terminal completion
required in Definition~\ref{def:law}. Proposition~\ref{cal:driver},
together with Proposition~\ref{eng:energy}, identifies its driver as a
weak solution of Equation~\eqref{mod:SDE} with
Equation~\eqref{mod:finite-energy}.

Apply Lemma~\ref{uni:weak} to the extension in
Lemma~\ref{uni:measurable}, with
$\kappa=16/3$ and $b=(2\pi/3)C$. Every admissible continuum
solution has the same driving law. A continuous driving path
determines its Loewner hulls. Whenever it has a trace, that trace at
each finite time is the boundary limit
\[
 \lim_{y\downarrow0}\phi^{-1}
             \bigl(g_t^{-1}(W_t+\ii y)\bigr),
\]
read in disc coordinates if a physical boundary value is ambiguous.
Thus the full capacity trace, its orientation and its continuous
terminal completion are determined by the driver. Inserting an
extra curve in a swallowed component would violate the trace
requirement and does not create another admissible curve law.

Existence requires only one bounded approximation. Such
approximations exist, for example on square rhombic meshes: exhaust
$\Omega$ by simply connected polygonal interiors, approximate a
chosen interior access sequence for each marked prime end, and
then refine the mesh diagonally. The polygonal approximation can
be kept in a fixed neighborhood of $\Omega$; compact kernel
convergence and convergence of the two prime-end coordinates are
preserved by the diagonal choice. With any nomes having the
required asymptotic ratio, compactness supplies a subsequential
limit and hence an admissible continuum law.

This use of a subsequence is an existence proof for the already
specified continuum equation; the definition and uniqueness do not
depend on its choice. For every bounded admissible approximation,
all tight subsequences have that unique law. The subsequence
criterion gives full-sequence convergence. Section~\ref{sec:localization}
removes the common diameter restriction.
\end{proof}

\begin{remark}
The uniqueness argument localizes energy on finite capacity horizons.
It needs neither a global Novikov condition nor a comparison of
complete infinite-horizon curve laws. In particular, no assertion
of mutual absolute continuity at infinite capacity is used in
Theorem~\ref{thm:main}.
\end{remark}
